\documentclass[10pt]{article}
\usepackage[margin=20mm]{geometry}
\usepackage{amsmath,amssymb,mathtools,fontspec,unicode-math,xurl,hyperref}
\IfFontExistsTF{DejaVu Serif}{\setmainfont{DejaVu Serif}}{\setmainfont{Latin Modern Roman}}
\setmathfont{latinmodern-math.otf}
\setlength{\parindent}{0pt}\setlength{\parskip}{6pt}\sloppy
\title{Scientific initial review of M25-0017}
\author{}
\date{}
\begin{document}
\maketitle
Abstract: The reviewed author-hosted proof develops Euclidean and hyperbolic characterizations of Weil–Petersson curves. All 78 source pages were read and visually inspected. Three core proof estimates fail in explicit admissible examples: the upper Fubini estimate for Möbius energy, the Euclidean distortion estimate for normal reflection, and the boundary-curvature estimate in the renormalized-area argument. Two peripheral claims require restricted parameter ranges. The first core estimate has an analytic repair supplied here; general repairs for the latter two remain unverified. A single-basepoint dyadic-cylinder interpretation and a smooth-welding interpretation are separately clarified, with the manuscript’s own counterevidence retained. No counterexample to the main Weil–Petersson equivalences is established.\par
\section*{Identity, version, and page convention}
Christopher J. Bishop, Weil-Petersson curves, \(\beta\)-numbers, and minimal surfaces, associated journal DOI 10.4007/annals.2025.202.1.2 (https://doi.org/10.4007/annals.2025.202.1.2). The actual input is the author-hosted black-and-white Annals proof at https://www.math.stonybrook.edu/\textasciitilde{}bishop/papers/wp\_annals\_bw.pdf. Its PDF metadata records creation on 17 June 2025. It has provisional pages 1–78, whereas the publisher metadata gives journal pages 111–188. The DOI is an identity link, not evidence of byte-for-byte or mathematical equivalence to the journal final.\par
All page references below mean the printed provisional page, equal to the one-based PDF page. The bound input is `papers/0017/0017.pdf`, 1,822,894 bytes, with SHA256\par
> 91d7444a417f577b2e154fa0f98ed05a3d1670fbf6c704f6730e367736d2b0f2.\par
The manifest SHA256 is\par
> 210893dcd1f1de5b6a62f7fc52f92126c4ffd0fd63016533dbcba50a720a545b.\par
The seven manifest inputs were rehashed successfully. No matching TeX source or paper-specific executable release was acquired; the proof PDF is the scientific input. Findings are version-bound and must be compared with the journal final before being described as final-publication defects.\par
\section*{Checked claims and outcome}
The review traced Theorems 1.1–1.6, Corollaries 1.7–1.8, Definitions 1–20, and the implication diagram on page 31. It read Sections 1–23 and Appendix A, including the references. The source text was read through all 7,821 extracted lines, with overlapping rereading where a displayed chunk was truncated. Every source page was also rendered and viewed; rendering alone was not counted as inspection.\par
The strongest conclusion is Major revision of the acquired version’s proof chain: Problems 1–3 exhibit invalid estimates used in central implications. Problem 1 is repaired analytically without changing its intended conclusion. Problems 2–3 require geometric estimates not completed here. Problems 4–5 are Minor range errors whose stated downstream applications avoid the counterexamples. Two further scope questions are left as Needs Evidence, rather than converted into unsupported counterexamples.\par
Source pages | Claims actually checked | Result and limitation\par
1–16 | Main statements, WP definitions, weldings, arclength and fractional Sobolev conditions | Quantifiers traced; smooth-welding ambiguity isolated in Section 5 (\#sec:clarifications). External function-theoretic equivalences were not reproved in full.\par
17–31 | Excess length, \(\beta\) numbers, Menger curvature, reflection, fat balls, convex hulls, curvature and area definitions | Normalizations and dependence diagram checked; dyadic quantifier and fat-ball definition clarified.\par
32–38 | Tangent/Whitney estimates; Möbius energy and excess length | Problem 1; correct domain repair given. Superadditive excess-length summations checked.\par
38–48 | Strengthened traveling salesman input; dyadic cylinders; Menger and reflection/fat-ball arguments | Dyadic series and ancestor summation checked. The full external traveling salesman and Pajot proofs were not independently reproduced.\par
49–55 | Hyperbolic distance, thickness, smoothed surfaces, normal reflection, minimal disks and renormalized area | Problems 2–4; all-minimal-surface scope left pending.\par
56–61 | Exhaustions, reverse area implication and higher-dimensional extension | Dependency and limiting arguments inspected; two unresolved geometric transitions recorded.\par
61–78 | Other energies, open questions, six additional characterizations and bibliography | Problem 5; appendix statements treated as surveyed dependencies, not fresh independent proofs.\par
\section*{Definitions, assumptions, and dependency limits}
A chord-arc Jordan curve \(\Gamma\) has a constant \(C\) such that the length \(\ell(x,y)\) of its shorter connecting subarc satisfies \(|x-y|\leq\ell(x,y)\leq C|x-y|\). Its unit tangent \(\tau\) exists almost everywhere. The central Euclidean comparison concerns \[\operatorname{M\ddot ob}(\Gamma)
 =\int_\Gamma\!\int_\Gamma\left(\frac1{|z-w|^2}-\frac1{\ell(z,w)^2}\right)|dz|\,|dw|
\quad\hbox{and}\quad
 \int_\Gamma\!\int_\Gamma\frac{|\tau(x)-\tau(y)|^2}{|x-y|^2}|dx|\,|dy|.\] These are nonnegative integrals; Tonelli is available even before finiteness is known. The admissible integration domain after interchanging variables still matters.\par
The geometric conditions include a scale-normalized \(\beta\) number for \(3Q\cap\Gamma\), a reflection-distortion number \(\rho(Q)\) measuring Euclidean departure from isometry, and a fat-ball quantity \(\varepsilon_\Gamma(Q)\). In hyperbolic upper-half-space the operative metric for the curvature/area calculations is \(ds_\rho^2=(|dx|^2+dt^2)/t^2\), of curvature \(-1\). The thickness \(\delta\) of the hyperbolic convex hull vanishes for a round circle. A totally geodesic hemisphere then has both principal curvatures zero. These elementary special cases obey, rather than evade, the hypotheses used in Problems 2 and 3.\par
For a minimal surface \(S\subset\mathbb H^3\), its principal curvatures are \(\pm\kappa\), and \(K=-1-\kappa^2\). For a regular height truncation \(S_t\), \[A_\rho(S_t)-L_\rho(\partial S_t)
 =-2\pi\chi(S_t)-\int_{S_t}\kappa^2\,dA_\rho
   +\int_{\partial S_t}(\kappa_g-1)\,dL_\rho.\] The final boundary integral must tend to zero to obtain the stated formula for renormalized area. Finiteness of an area integral by itself is not a pointwise or trace estimate.\par
The paper’s existence assertions about a minimal disk, its assertions about every minimal surface with a given boundary, and its statements for a surface with finite Euler characteristic are different quantifiers. They are kept separate throughout this report. None of the elementary counterexamples below refutes the existence of a suitable minimal spanning surface for a WP curve.\par
\section*{Supported substantive problems}
Problem 1 (Major: the upper Fubini estimate uses a divergent enlargement).  Location and source claim. Lemma 9.1, especially the last chain of comparisons on page 36, is used to identify Definitions 7 and 8 and hence to prove Theorem 1.2. With \(s=\ell(\Gamma)/2\), the printed upper comparison replaces the domain \(\sigma(x,y)=\{(z,w):x,y\in\gamma(z,w)\}\) by the product of two length-\(s\) arcs and writes \[J(x,y)=
 \int_{z\in\sigma(x,y,s)}\int_{w\in\sigma(y,x,s)}
       |x-w|^{-4}|dz|\,|dw|
 \simeq
 \int_{w\in\sigma(y,x,s)}|x-w|^{-3}|dw|
 \simeq |x-y|^{-2}.\] Here \(\gamma(z,w)\) is the shorter arc. The first displayed comparison is false.\par
Complete admissible check. Take the unit circle, a smooth closed chord-arc Jordan curve, with angular coordinates \(x=0\) and \(y=d\), \(0<d<\pi/8\). Its total length is \(2\pi\), so \(s=\pi\). The two arcs in the printed product are \([-\pi,0]\) and \([d,d+\pi]\), understood modulo \(2\pi\). Since the integrand is independent of \(z\), \[J(d)=\pi\int_d^{d+\pi}\frac{dw}{(2\sin(w/2))^4}
      \sim \frac{\pi}{3d^3}\qquad(d\downarrow0).\] The last claimed quantity is \((2\sin(d/2))^{-2}\sim d^{-2}\), so no uniform comparison holds. On this circle \(|\tau(x)-\tau(y)|^2=|x-y|^2\); the final tangent integral is finite, whereas the printed intermediate integral has a local logarithmic divergence after multiplication by this tangent difference and integration in \(d\).\par
Merely changing the denominator to \(|z-w|^4\) does not repair the argument. The two semicircles overlap on an arc of length \(d\) (represent the first as \([\pi,2\pi]\)). Their product contains an open segment of the diagonal, and \(\int\int|z-w|^{-4}|dz|\,|dw|=\infty\). Thus a domain correction is necessary in addition to any typographical correction.\par
Analytic repair. Retain the shorter-arc condition. For \(d=\ell(x,y)\leq\ell(\Gamma)/8\), parametrize \(z\) at backward arclength \(a\geq0\) from \(x\) and \(w\) at forward arclength \(b\geq0\) from \(y\). The correct domain satisfies \[a+b+d\leq \ell(\Gamma)/2.\] The chord-arc hypothesis gives \(|z-w|\geq C^{-1}(a+b+d)\) on this domain. Including the opposite orientation changes only a factor of two. Consequently the inner kernel is bounded above by \[2C^4\int_0^\infty\int_0^\infty(a+b+d)^{-4}\,da\,db
 =\frac{C^4}{3d^2}.\] For a lower bound, restrict to \(0\leq a,b\leq d\). This is admissible because \(3d\leq3\ell(\Gamma)/8<\ell(\Gamma)/2\), and the chord length is at most \(a+b+d\leq3d\). The kernel is at least \(d^2/(3d)^4=1/(81d^2)\) for one orientation. Pairs \((x,y)\) outside this small-separation region contribute a finite bounded amount: every containing shorter arc has length at least \(\ell(\Gamma)/8\), hence its chord is bounded below by the chord-arc hypothesis. This proves the intended comparability using the actual domain.\par
Counterevidence, impact, and repair status. The circle satisfies both claimed finiteness conclusions. The preceding proof repairs the central comparison and supports the theorem rather than refuting it. Nevertheless the printed upper derivation is invalid in an admissible, maximally regular case and lies directly on the proof of Theorem 1.2, which supports Major under standard 0.7.1. The exact quadrant integral and illustrative values of \(d^3J(d)\) are recorded in `independent\_checks.json`; the analytic proof above, not those finite numerical samples, establishes the repair.\par
Problem 2 (Major: hyperbolic flatness does not eliminate Euclidean reflection distortion).  Location and source claim. In the proof of Lemma 18.2 on page 53, the center of the boundary sphere of the normal hyperplane at \(z=(x,t)\) is asserted to lie within \(O(t\sup_w\max_j|\kappa_j(w)|)\) of \(x\). The normal antipodal construction is then used to assert \(\rho(Q)=O(K_r(z))\), where \(K_r\) bounds the principal curvatures on a fixed hyperbolic ball. This is the quantitative step in the higher-dimensional implication (16)\(\Rightarrow\)(13).\par
Parameters and assumptions. Set \(n=2\), let \(\Gamma=\{|z|=1\}\subset\mathbb R^2\), and let \[S=\{(x,y,t):x^2+y^2+t^2=1,\ t>0\}\subset\mathbb H^3.\] This is a smooth, properly embedded, complete, totally geodesic disk with asymptotic boundary \(\Gamma\). Both principal curvatures are zero everywhere. In particular, they tend to zero at the boundary and their square integral is zero, so the surface obeys the curvature requirements of Definition 16. Any fixed-ball supremum \(K_r\) is zero.\par
At \(p=(x,0,t)\in S\), \(0<x<1\), the Euclidean center of the normal geodesic is at horizontal coordinate \(a=1/x\). Its endpoints are \[r_-=\frac{1-t}{x},\qquad r_+=\frac{1+t}{x},
 \qquad r_-r_+=1.\] These formulas follow by taking a Euclidean circle orthogonal to the plane \(t=0\) and tangent at \(p\) to the Euclidean normal \((x,0,t)\); the upper-half-space metric is conformal, so the normal direction is unchanged. Thus \[a-x=\frac{t^2}{x}>0,\] whereas the asserted \(O(t\sup|\kappa|)\) is identically zero. This remains a boundary counterexample: take \(x=(1-u^2)/(1+u^2)\), \(t=2u/(1+u^2)\) and let \(u\downarrow0\).\par
The antipodal pairing of the two endpoints is precisely \[R(z)=\frac1{\overline z}.\] Its Euclidean differential has norm \(|z|^{-2}\). In a nondegenerate Whitney cube off the unit circle this differs from one at some point, so the local distance distortion \(\rho(Q)\) is positive. The proposed bound by \(K_r=0\) is false.\par
Alternative readings and counterevidence. Hyperbolic isometric normalization at a point cannot remove the Euclidean scale distortion in the definition of \(\rho(Q)\); a Möbius transformation is not a Euclidean isometry. For this same circle, inversion does satisfy \(\rho(Q)=O(\operatorname{diam}Q)\) near the circle, and the resulting square sum is finite: there are \(O(h^{-1})\) relevant cubes at side length \(h\), each contributing \(O(h^2)\). Hence the existence of a reflection in Definition 13 survives.\par
The paper itself recognizes on page 49 that \(\delta=0\) for a circle while \(\varepsilon_\Gamma\) is nonzero, so a size term is needed when passing from hyperbolic to Euclidean geometry. That observation is relevant counterevidence against reading the omitted scale term as harmless.\par
Impact and remaining repair. The displayed estimate cannot prove the central implication as written. A possible form of a replacement must include a scale or background-circle term, for example a properly normalized \(O(\operatorname{diam}Q)\) contribution, in addition to curvature information. This report does not assert a general bound of that form without proof. One must also justify replacing an \(L^2\) curvature integral by an \(L^2\) integral of fixed-ball suprema. For a general smooth function such a replacement fails: disjoint bumps of amplitude \(k^{-1/2}\) supported in disks of radius \(2^{-k}\) have finite square integral, whereas their fixed-radius local suprema have divergent square integral on disjoint fixed-radius neighborhoods. This diagnostic is not a constructed surface counterexample, but identifies an additional estimate the repair must supply. No counterexample to the existential equivalence is claimed.\par
Problem 3 (Major: the boundary error estimate fails on a totally geodesic disk).  Location and source claim. Lemma 20.1, page 55, asserts \(\kappa_g=1+O(\delta^2)\) and equation (20.1), \[\int_{\partial S_t}(\kappa_g-1)\,dL_\rho
       =O\!\left(\int_{\partial S_t}\delta^2\,dL_\rho\right),\] then uses a trace-to-area estimate to make the left side vanish. This is central to (15)\(\Rightarrow\)(19), Corollary 1.7 and formula (1.6).\par
Full check. Use the same unit hemisphere as in Problem Problem 2 (\#prob:reflection). Its round boundary is WP, its convex hull is the totally geodesic plane itself, \(\delta\equiv0\), \(\kappa\equiv0\), and \(\chi(S)=1\). Every \(0<t<1\) is a regular truncation height, so all conditions of the displayed calculation hold.\par
Write \(r=\sqrt{1-t^2}\). Direct integration of the induced hyperbolic metric gives \[L_\rho(\partial S_t)=\frac{2\pi r}{t},\qquad
 A_\rho(S_t)=2\pi\left(\frac1t-1\right).\] Intrinsically \(S\) is the hyperbolic plane and this section is a geodesic circle with \(\cosh R=1/t\). Therefore \[\kappa_g=\coth R=\frac1{\sqrt{1-t^2}}>1,\qquad
 \int_{\partial S_t}(\kappa_g-1)\,dL_\rho
   =\frac{2\pi(1-\sqrt{1-t^2})}{t}
   =\pi t+O(t^3)>0.\] The right side of (20.1) is zero. The assertion that the horizontal curvature contribution is bounded by the surface’s principal curvatures also fails in this example: the surface is totally geodesic, while its horizontal sections are Euclidean circles.\par
Counterevidence and exact consequence. The error nevertheless tends to zero. Indeed \[\lim_{t\downarrow0}\{A_\rho(S_t)-L_\rho(\partial S_t)\}=-2\pi,\] which agrees with the claimed renormalized-area formula. Thus this is a counterexample to the proof estimate, not to Corollary 1.7. The example exhibits the missing background term explicitly: \(\kappa_g-1=\tfrac12t^2+O(t^4)\) even when \(\delta=0\).\par
Repair requirements. A valid argument must retain such a scale term and control its integrated boundary contribution for the surfaces under consideration, as well as establish a genuine trace estimate. On the round disk an \(O(t^2)\) term integrates to \(O(t)\), but this special case is not a proof of the general estimate. The printed annulus \(A_t=S_t\setminus S_{t+1}\) does not escape compact subsets as \(t\downarrow0\), so finite total integral does not imply its integral tends to zero. A geometric-height annulus such as \(S_t\setminus S_{2t}\) would have the correct escaping property; its required comparison with the boundary trace still needs proof. The formula involving distance from an interior point at the end of page 55 concerns geodesic spheres and retains a \(\coth r\) term; it does not justify dropping the corresponding error for height sections. General repair and the all-minimal-surface scope remain unverified.\par
Problem 4 (Minor: the half-space distance comparison needs a small-gap hypothesis).  Location and claim. Lemma 16.1, page 49, states that disjoint Euclidean balls of equal radius \(r\) separated by \(\varepsilon\) have corresponding hyperbolic half-spaces at distance \(\rho\simeq\sqrt{\varepsilon/r}\), without restricting \(\varepsilon/r\).\par
Complete calculation. Let \(\eta=\varepsilon/r>0\). In the two-dimensional section through the centers the boundary geodesics are semicircles of radius \(r\) with centers \(2r+\varepsilon\) apart. Their inversive distance gives \[\cosh\rho=\frac{(2r+\varepsilon)^2-2r^2}{2r^2}
          =1+2\eta+\frac{\eta^2}{2},
 \qquad
 \rho=2\operatorname{arcosh}(1+\eta/2).\] These balls are disjoint for every \(\eta>0\), so they meet all printed assumptions. At small \(\eta\), \(\rho/\sqrt\eta\to2\). At large \(\eta\), \(\rho=2\log\eta+O(\eta^{-1})\) and \(\rho/\sqrt\eta\to0\). The stated two-sided uniform comparison therefore fails.\par
Counterevidence, impact and verified correction. The proof uses a local expansion whose optimizing height would exceed the radius when the gap is large. The next lemma’s application explicitly requires \(\varepsilon_\Gamma(Q)<1/100\), so the relevant small-gap use is not refuted. Restricting \(0<\eta\leq\eta_0\) for a fixed finite \(\eta_0\) gives a valid square-root comparison, by continuity and the limit at zero. The exact formula above works for all positive gaps. This is a Minor range error, not a failure of the subsequent small-thickness implication.\par
Problem 5 (Minor: smooth curves need not have finite knot energy at the claimed endpoint).  Location and claim. Section 23, page 61, defines \[E^{j,p}(\Gamma)
 =\int_\Gamma\!\int_\Gamma
   \left(|x-y|^{-j}-\ell(x,y)^{-j}\right)^p\,|dx|\,|dy|\] and asserts finiteness for smooth curves if \(jp\leq2p+1\).\par
Counterexample with all conditions. Take the smooth embedded unit circle, \(j=3\) and \(p=1\). These parameters satisfy \(jp=2p+1=3\). For shorter arclength \(h\downarrow0\) the chord is \(2\sin(h/2)\), and \[(2\sin(h/2))^{-3}-h^{-3}
 =\frac1{8h}+\frac{17h}{1920}+O(h^3).\] The positive integrand consequently has a logarithmically divergent integral near the diagonal. Thus \(E^{3,1}(\Gamma)=\infty\) despite the asserted smoothness and endpoint condition.\par
Repair and limits. For a unit-speed smooth curve with nonzero curvature at a point, \[|x(s+h)-x(s)|=|h|\left(1-\frac{\kappa(s)^2h^2}{24}+O(h^3)\right).\] For \(j,p>0\) the energy integrand has leading order \((j\kappa(s)^2/24)^p|h|^{p(2-j)}\). Local integrability at nonzero curvature requires the strict inequality \(jp<2p+1\). The unit circle calculation rigorously excludes the printed equality case without relying on any numerical integration. The main Möbius energy corresponds to \(j=2,p=1\) and is strictly inside this range. The issue is therefore peripheral and Minor; no failure of Theorem 1.2 is inferred from it.\par
\section*{Scope clarifications and counterevidence}
These items are not added to the supported Major/Minor totals. Their purpose is to state exactly which reading is justified and to prevent an overbroad adverse conclusion.\par
\subsection*{A single dyadic cylinder is insufficient, but the main theorem already quantifies correctly}
Definition 20 on page 28 refers to the dyadic cylinder \(X\) corresponding to \(\Gamma\). Read in isolation as a condition on one basepoint, it is too weak. Let \(\Gamma\) be the boundary of the unit square, of length four, and choose a vertex as the arclength origin. Its \(2^n\)-vertex equalarclength polygon equals the square for \(n\geq2\), with only collinear subdivision. For \(n=1\) its closed two-edge length is \(2\sqrt2\), so \[\sum_{n\geq1}2^n\{\ell(\Gamma)-\ell(\Gamma_n)\}=8-4\sqrt2<\infty.\] Using exactly the page-28 layers \(\Gamma_n\times(2^{-n},2^{-n+1}]\), for \(0<t<1/2\) the lateral area and bottom-section length satisfy \[A_\rho(X_t)=2\sqrt2+4(1/t-2),\qquad
 L_\rho(\partial X_t)=4/t.\] The renormalized area is the finite constant \(2\sqrt2-8\). Adding or changing finitely many upper layers or caps only changes a finite constant. The harmless indexing inconsistency in Section 12 does not affect this example.\par
The square is not WP. Across a corner, the two unit tangents differ by a fixed nonzero vector and \[\int_0^\varepsilon\int_0^\varepsilon\frac{da\,db}{(a+b)^2}=\infty.\]\par
Equivalently, for dyadic \(Q=(-h,0]^2\) adjacent to a corner at the origin, \(3Q\cap\Gamma\) contains \((0,0),(h,0),(0,h)\). The smallest strip half-width containing this triangle is \(h/(2\sqrt2)\). With the paper’s normalized \(\beta\) this gives \(\beta_\Gamma(Q)\geq1/(3\sqrt2)\) at every sufficiently fine scale, and the square sum diverges.\par
The omitted uniformity is real: shifting the arclength origin by \(1/3\) yields corner fractions \(1/3\) and \(2/3\) at every dyadic scale \(n\geq2\), so \[2^n\{\ell(\Gamma)-\ell(\Gamma_n)\}
 =16\left(1-\frac{\sqrt5}{3}\right)>0.\] This series diverges.\par
Counterevidence controlling the judgment. Theorem 1.3 on page 4 and Lemma 12.1 on page 40 explicitly require every basepoint with a uniform bound. Those statements exclude this square and are not refuted. Definition 20’s reference to Theorem 1.3 can reasonably be read as importing that qualification. Accordingly this report records a scope clarification: carry the universal basepoint and uniform-bound requirement explicitly into Definition 20 and its use in Lemma 22.1 on page 60. It does not count the abbreviated wording as a separate demonstrated core defect.\par
\subsection*{Smooth weldings require a nondegeneracy qualification}
The transition across pages 15–16 says that an \(H^s\) welding and an \(H^s\) arclength parametrization are equivalent for \(s>3/2\), using Hölder continuity of the welding derivative to infer nonvanishing conformal derivatives and biLipschitz behavior. Hölder continuity alone does not imply a positive lower bound.\par
For the unqualified interpretation there is an explicit test. Let \(M(z)=z/(1+z)\), \[f(z)=M(z^{1/2})\quad(z\in\mathbb H^+),\qquad
 g(w)=M(w^{3/2})\quad(w\in\mathbb H^-),\] with arguments in \((0,\pi)\) and \((-\pi,0)\), respectively. The image of \(f\) is the upper half-disk bounded by the segment \([0,1]\) and the upper semicircle of radius \(1/2\); the image of \(g\) is its complementary domain on the sphere. Both are conformal bijections. Their welding on the extended real line is \(f^{-1}g(t)=t^3\): for \(t<0\), the lower-branch power is \(i|t|^{3/2}\) and the upper-branch square root of \(-|t|^3\) is the same. After Cayley changes of source coordinates the welding is a smooth circle homeomorphism. At infinity its local expression is again \(u^3\), and it is quasisymmetric. It belongs to every \(H^s\), but its derivative vanishes at two points. The curve has two right-angle corners, so its arclength tangent is not in \(H^{1/2}\) by the preceding corner integral.\par
The next paragraph explicitly discusses \(H^s\) diffeomorphisms, and an implicit restriction to the WP class could also exclude this example. Consequently the precise correction is to distinguish an \(H^s\) circle homeomorphism from a nondegenerate \(H^s\) diffeomorphism and state the intended domain of the equivalence. This report does not count the example as refuting a theorem restricted to WP weldings or diffeomorphisms. It also does not claim to have reproved the full higher-\(s\) equivalence under the corrected hypotheses.\par
\subsection*{Restore the scale and incidence conditions in the fat-ball definition}
The planar wording around Definition 14 on page 23 must be read with the incidence and normalization visible in Figure 6 and the later construction: the near-contact point lies in the relevant cube and the gap is measured relative to its diameter, while the ball radii are controlled from below at the appropriate inverse scale. If one omits incidence and uses only a line through \(Q\), two arbitrarily large balls can be placed arbitrarily far along that line, separated by an arbitrarily small gap and missing any fixed bounded curve. The resulting infimum would vanish even for a square. The later proof on pages 46–48 uses the intended local geometry. This is recorded as restoration of the definition, not inflated into a separate scientific error after ignoring its figure and usage.\par
Other observed notation or extraction discrepancies are excluded from the issue count. In particular, the page-18 image includes the square root in the triangle-curvature expression, and the page-7 curvature identity has the intended absolute-value bars. Fractional tangent phases require their circle factor and orientation convention; these do not produce a new counterexample to Theorem 1.1. Obvious indexing, naming, and survey-formula slips are not used as substitutes for substantive findings.\par
\section*{Needs Evidence and dependence of the overall claims}
\subsection*{NE1: the quantifier over all minimal surfaces}
Theorem 1.6(4), page 7, and page 55 involve every minimal surface, whereas Lemma 19.1 on page 53 is stated for an embedded minimal disk. The step from small convex-hull thickness to uniform local curvature, a single boundary sheet, finite Euler characteristic, and the integrated estimate requires hypotheses that control this distinction.\par
Selected portions of Seppi’s arXiv:1411.3412v4 were inspected: Lemma 4.11, the local compactness discussion, Proposition 4.14, Lemma 4.15 and formula (32). Proposition 4.14 allows a fixed quasicircle bound \(K_0\) and a fixed radius, with constants depending on them; it must not be misrepresented as requiring \(K_0\) arbitrarily close to one merely to state that local estimate.\par
Huang–Lowe–Seppi, arXiv:2309.00599v2, Section 1.4, explicitly discusses a stability requirement in the earlier compactness argument, and removes it in a globally small-quasicircle regime using a new uniqueness result. This is a relevant later dependency, not evidence that every arbitrary WP boundary satisfies that global smallness hypothesis. Only the indicated selected portions were consulted, not either entire external proof.\par
The pending task is to give a complete localization and stability/topology argument for every surface in the original quantifier, or explicitly restrict the affected intermediate result. No surface violating the final theorem with a WP asymptotic boundary has been constructed here. NE1 therefore remains Needs Evidence rather than a supported theorem counterexample.\par
\subsection*{NE2: the reverse exhaustion argument and boundary regularity}
Definition 18 on page 28 assumes finite Euler characteristic and a nested compact exhaustion with controlled length minus area. In Section 21, pages 56–59, the proof uses a height end that is an annulus, a small-curvature regime, and then a sharper Euler-characteristic constant. Finite Euler characteristic alone does not identify a particular height function’s regular sublevel sets with the desired annular exhaustion. The implication to the stated height-based limit, especially in the non-WP case where the desired value is \(-\infty\), needs those transitions supplied.\par
Similarly, the extension of the boundary graphical and normal-foliation arguments from the regularity assumptions of the cited Lin result to the WP setting is asserted in Sections 18 and 22. The full Lin paper was not acquired and checked here. The report records this dependency rather than inventing a regularity theorem. The Fiala-type differential inequality and its intended geometric use were inspected; a complete proof of all the exhaustion and foliation prerequisites was not obtained. This is a remaining proof obligation, not a fabricated counterexample.\par
\subsection*{What survives, what needs repair, and later versions}
Problems 1–3 affect different edges of the equivalence diagram. The correct Fubini argument restores the tangent/Möbius edge. The normal-reflection and height-trace bounds need additional geometric work; their round examples show exactly why an estimate depending only on extrinsic hyperbolic curvature or thickness is insufficient. These facts do not negate the other Euclidean characterizations, nor do they show that the stated renormalized-area limit is false. Problems 4 and 5 have restricted, analytically verified repairs and do not undermine their cited main applications.\par
A narrow public search for an author correction or an erratum specific to this paper did not identify one. This is a bounded negative search, not proof that no correction exists. The author-hosted proof and the journal final were not compared, so adoption or correction status in the final publication is unknown. The selected later minimal-surface paper is used only to clarify a dependency. No new empirical experiment or author-code execution is needed to verify the explicit failures in this report; general geometric repair requires proof.\par
\section*{Conclusion and remaining work}
Item | Judgment | Repair verification | Adoption / remaining work\par
P1 | Major | Correct shorter-arc domain gives both kernel bounds analytically | Author adoption unrequested; compare final text.\par
P2 | Major | Round counterexample verified; general scale-corrected estimate not proved | Supply distortion, foliation and local-supremum estimates.\par
P3 | Major | Round counterexample verified; general trace repair not proved | Supply height-section error and trace control.\par
P4 | Minor | Exact distance formula; small-gap range verified | Restrict statement or use exact formula.\par
P5 | Minor | Circle endpoint divergence; strict local threshold verified | Exclude equality in the general smooth-curve claim.\par
S1–S3 | Scope clarification | Explicit square and welding tests; contextual counterevidence retained | Make quantifiers, nondegeneracy and normalization explicit.\par
NE1–NE2 | Needs Evidence | No theorem counterexample asserted | Complete stability, topology, boundary and exhaustion arguments.\par
This is an initial scientific review, not independent acceptance. Human scientific approval and author adoption have not been supplied. The concrete outcome is three supported Major proof defects, two supported Minor local claims, two Needs Evidence items, and three scope clarifications.\par
9 C. J. Bishop, Weil-Petersson curves, \(\beta\)-numbers, and minimal surfaces, author-hosted black-and-white Annals proof, PDF created 17 June 2025, provisional pp. 1–78, bound SHA256 in Section 1. https://www.math.stonybrook.edu/\textasciitilde{}bishop/papers/wp\_annals\_bw.pdf. Associated DOI: https://doi.org/10.4007/annals.2025.202.1.2.\par
A. Seppi, Minimal discs in hyperbolic space bounded by a quasicircle at infinity, arXiv:1411.3412v4, 9 November 2016, 24-page author version. https://arxiv.org/abs/1411.3412v4. Selected portions around Lemma 4.11, Proposition 4.14, Lemma 4.15 and formula (32) consulted; not a full independent review. Journal identifier: https://doi.org/10.4171/CMH/403.\par
Z. Huang, B. Lowe and A. Seppi, Uniqueness and non-uniqueness for the asymptotic Plateau problem in hyperbolic space, arXiv:2309.00599v2, 17 September 2024. https://arxiv.org/abs/2309.00599v2. Introduction and Section 1.4 selected passages consulted for scope; whole-paper verification not claimed.\par
\section*{Supporting checks}
It imports installed SymPy and mpmath, performs the nine check groups described below, and writes the adjoining JSON. It neither imports nor executes author code and uses no network. The actual command was the existing Windows Python executable with `-X utf8 -B`, followed by that full staging-local script path. Versions were Python 3.13.5, SymPy 1.13.3 and mpmath 1.3.0.\par
The nine groups check the aligned-square area identity; the shifted-square constant deficit; the exact double kernel integral; numerical illustrations of the erroneous circle kernel; the hemisphere’s normal endpoints and center; the hemisphere boundary trace and area limit; the power-welding boundary identities; the exact half-space separation limits; and the circle knot-energy series. All recorded assertions passed. The floating-point values are explicitly illustrative. For example \(d^3J(d)\) for \(d=10^{-1},10^{-2},10^{-3}\) is approximately \(1.0524484,1.0472499,1.0471981\), approaching \(\pi/3\); the divergent comparison is proved by the analytic asymptotic, not inferred from these three values.\par
The reviewer-owned extraction helper used existing Poppler to generate `source\_pages.json` and the complete page-delimited text. The existing rendering helper produced all 78 PNGs. Every PNG was opened for visual inspection. Pages 65–68 were explicitly reopened after one earlier tool display had been truncated; the reopened images, not the truncated display, support their inspection status.\par
\section*{Supporting checks}
The complete dispatch read receipt, metadata amendment and hash/capability ledger are stored at the staging root’s `evidence/` directory. The receipt records the full packet, original task, scientific prompt and all ten manifests read to zero remaining lines before scientific review. The source-preparation acquisition token is different and is not substituted for this review token.\par
The progressive `delivery\_state.json` and per-paper `delivery\_inventory.json` contain actual file sizes and hashes. The latter is the authoritative output inventory after finalization; a self-referential report hash is not fabricated inside this report.\par
No installation was attempted. This report was not compiled; no file at `reports/Pdf/0017.pdf` is claimed. Static environment-balance checks do not establish compilability or visual quality. The source-paper PNGs are not a substitute for a compiled report or report-page inspection.\par
The requesting user authorized the task and its local Windows destination. All scientific reasoning, reviewer calculations, source-page inspection and report drafting were performed by this assistant. No scientific subagent or human review is claimed. These are separately sourced labels, not an asserted equivalence or backend-route verification. No independent evidence of a non-6-Pro route was observed, and no browser/model switch or denied-tab takeover was attempted.\par
All written artifacts remain within the task’s authorized staging directory. There was no final-archive write, author contact, author-code execution, software installation, payment, quota request, permission expansion, ZIP/base64 export, Library/cloud upload or heartbeat creation. Independent acceptance remains pending.\par
\section*{Compilation continuation (26 October 2026)}
Scientific conclusions, pending questions, source-version limits and the absence of independent acceptance remain unchanged.\par
\end{document}
